\documentclass[11pt,a4paper]{article}

\usepackage[a4paper,margin=1in]{geometry}
\usepackage{fontspec}
\IfFontExistsTF{Times New Roman}
  {\setmainfont{Times New Roman}}
  {\IfFontExistsTF{TeX Gyre Termes}{\setmainfont{TeX Gyre Termes}}{\setmainfont{Times}}}
\usepackage{newtxmath}
\usepackage{amsmath}
\usepackage{booktabs}
\usepackage{longtable}
\usepackage{array}
\usepackage{enumitem}
\usepackage{xcolor}
\usepackage{listings}
\usepackage{graphicx}
\usepackage{float}
\usepackage[round,authoryear]{natbib}
\usepackage[hidelinks]{hyperref}
\usepackage{xurl}
\usepackage{fancyhdr}
\usepackage{microtype}

\setlength{\parindent}{0pt}
\setlength{\parskip}{0.65em}
\setlength{\emergencystretch}{3em}
\setlist[itemize]{leftmargin=1.4em,itemsep=0.2em,topsep=0.2em}
\setlist[enumerate]{leftmargin=1.6em,itemsep=0.2em,topsep=0.2em}

\pagestyle{fancy}
\fancyhf{}
\lhead{\small Ticket 09 Study Report}
\rhead{\small Continuation Premium and Boundary Extraction}
\cfoot{\thepage}
\renewcommand{\headrulewidth}{0.3pt}

\lstset{
  basicstyle=\ttfamily\small,
  breaklines=true,
  frame=single,
  columns=fullflexible,
  keepspaces=true,
  rulecolor=\color{black!30},
  xleftmargin=0.5em,
  xrightmargin=0.5em
}

\newcommand{\code}[1]{\texttt{#1}}
\newcommand{\file}[1]{\path{#1}}
\newcommand{\func}[1]{\path{#1}}
\newcolumntype{L}[1]{>{\raggedright\arraybackslash}p{#1}}
\newcolumntype{P}[1]{>{\raggedright\arraybackslash\ttfamily\footnotesize}p{#1}}

\title{\textbf{Ticket 09 Study Report: Continuation Premium and Boundary Extraction}\\
\large SURF2026 American Option Risk Surfaces}
\author{Codex-assisted implementation report}
\date{June 16, 2026}

\begin{document}
\maketitle

\begin{abstract}
Ticket 09 is the first project ticket where free-boundary extraction is intentionally in scope.
It uses existing American CN/PSOR outputs and computes the continuation premium
\(U-\Phi\), where \(U\) is the American option value and \(\Phi\) is the payoff. The extraction
layer then looks for option-type-aware transitions between exercise-like and continuation-like
regions. This report explains the method for beginner researchers who know Python and data work
but are still learning American option free-boundary diagnostics.
\end{abstract}

\tableofcontents
\newpage

\section{Role of Ticket 09 in the Whole Project}

Tickets 01 through 04 built the pricing machinery: payoff functions, Black-Scholes utilities,
finite-difference grids, European Crank--Nicolson validation, and the American CN/PSOR core.
Tickets 05 through 07 validated American put, no-dividend American call, and dividend-paying
American call behavior. Ticket 08 then checked obstacle and complementarity diagnostics without
extracting a boundary.

Ticket 09 is the natural next layer. It takes the stored American value grid and asks:
\begin{itemize}
  \item Where does the value equal payoff, up to a stated threshold?
  \item Where does the value become strictly greater than payoff?
  \item Can we record the transition as an approximate boundary without pretending it is exact?
\end{itemize}

The output is a boundary-extraction diagnostic, not a new pricing algorithm. The solver core is
unchanged.

\section{Theory and Methodology Connection}

The documented source basis for Ticket 09 is internal:
\begin{itemize}
  \item the planning report, which places boundary extraction after solver validation and
  complementarity diagnostics;
  \item the literature map, which emphasizes that boundary extraction should use continuation
  premium rather than visual guessing;
  \item the formulation note, especially the sections on continuation premium and free-boundary
  extraction;
  \item the solver validation plan, especially its boundary extraction validation guidance;
  \item the Ticket 04 through Ticket 08 implementation reports and tests.
\end{itemize}

This reference integration pass adds verified citations for the external methodology behind the
project documents. The approximate free-boundary and finite-difference context is connected to
Brennan-Schwartz and computational-finance PDE background
\citep{brennan_schwartz_1977,wilmott_howison_dewynne_1995}, while complementarity terminology is
connected to standard LCP methodology \citep{cottle_pang_stone_2009}.

\section{Continuation Premium \(U-\Phi\)}

For an American option, the value \(U(S,\tau)\) should not fall below the immediate exercise
payoff \(\Phi(S)\). Define the continuation premium:
\[
P_{\mathrm{cont}}(S,\tau) = U(S,\tau) - \Phi(S).
\]

In plain language, this is the extra value from holding the option instead of exercising
immediately. The interpretation is:
\begin{itemize}
  \item \(P_{\mathrm{cont}} \approx 0\): the value is close to payoff, so the node is
  exercise-like;
  \item \(P_{\mathrm{cont}} > 0\): waiting has extra value, so the node is continuation-like;
  \item the free boundary is approximated where the premium crosses from near zero to positive.
\end{itemize}

The word ``approximately'' matters. A finite-difference grid has spacing error, PSOR tolerance,
and a discrete threshold. Ticket 09 therefore reports the threshold and no-boundary reasons.

\section{Why Boundary Extraction Comes After Diagnostics}

Boundary extraction is more fragile than a price comparison. A price can look reasonable while
the obstacle or complementarity conditions fail. That is why Ticket 08 came first: it checked that
the representative American put and dividend-paying American call results satisfy the LCP
diagnostics at interior nodes.

Ticket 09 uses that groundwork. It does not re-solve the option and does not recompute LCP
residuals. It reads the stored value grid, computes \(U-\Phi\), and extracts approximate transition
points. If future diagnostics fail, boundary estimates should be treated with suspicion.

\section{Put Search Direction}

For an American put, early exercise is expected at low spot prices. When the stock is very low,
the put is deep in the money, and immediate exercise can be attractive. Therefore the search rule
is:
\begin{enumerate}
  \item move from low \(S\) toward high \(S\);
  \item identify the low-\(S\) exercise-like run where \(P_{\mathrm{cont}}\le\epsilon\);
  \item identify the next continuation-like run where \(P_{\mathrm{cont}}>\epsilon\);
  \item interpolate the threshold crossing between the last exercise-like node and first
  continuation-like node.
\end{enumerate}

The far high-\(S\) tail of a put can also have very small option value and therefore very small
premium. Ticket 09 does not treat that far tail as the exercise boundary. It uses the option-type
search direction to find the economically relevant low-spot transition.

\section{Dividend-Call Search Direction}

For a dividend-paying American call, early exercise can become meaningful at high spot prices.
Exercising early sacrifices optionality, but it may let the holder own the stock before dividends
reduce its value. Therefore the search rule is:
\begin{enumerate}
  \item move from high \(S\) toward lower \(S\);
  \item identify the high-\(S\) exercise-like run where \(P_{\mathrm{cont}}\le\epsilon\);
  \item identify the adjacent continuation-like run below it;
  \item interpolate the threshold crossing between the continuation-like and exercise-like nodes.
\end{enumerate}

The low-\(S\) tail of an out-of-the-money call can also have tiny value and tiny premium. Ticket
09 ignores that low-spot tail for call boundary extraction, because the dividend-call exercise
region of interest is the high-spot side.

\section{Threshold and Interpolation Choices}

Ticket 09 uses the default threshold
\[
\epsilon = 10^{-6}.
\]
The rule is:
\[
\text{exercise-like if } P_{\mathrm{cont}}(S_i,\tau) \le \epsilon,\qquad
\text{continuation-like if } P_{\mathrm{cont}}(S_i,\tau) > \epsilon.
\]

If a clean transition exists between adjacent nodes \(S_i\) and \(S_{i+1}\), Ticket 09 uses
linear interpolation:
\[
S_f
= S_i
+ \frac{\epsilon - P_i}{P_{i+1}-P_i}
\left(S_{i+1}-S_i\right),
\]
where \(P_i=P_{\mathrm{cont}}(S_i,\tau)\). This is a practical grid diagnostic, not an exact
analytical boundary formula.

Boundary extraction decisions use interior spot nodes to reduce imposed-boundary artifacts.
Full-grid premium profiles are still stored and plotted for study purposes.

\section{No-Boundary Metadata}

Ticket 09 does not silently return a missing value. Each time row records a reason when no
boundary is extracted:

\begin{longtable}{P{0.36\textwidth}L{0.50\textwidth}}
\toprule
Reason & Meaning \\
\midrule
\endhead
\path{maturity_row_ambiguous} & At \(\tau=0\), value equals payoff everywhere, so the whole row is exercise-like by construction. \\
\path{all_continuation_like} & No interior node has premium below the threshold. \\
\path{all_exercise_like} & Every interior node is threshold-exercise-like. \\
\path{expected_exercise_side_absent} & The option-type side where exercise is expected does not show threshold-exercise behavior. \\
\path{no_clean_transition} & The observed pattern is not a stable transition under the extraction rule. \\
\path{insufficient_interior_nodes} & There are not enough interior nodes to bracket a transition. \\
\bottomrule
\end{longtable}

\section{Continuation Premium Versus American-Minus-European Value}

Continuation premium is not the same as American-minus-European value.
\[
P_{\mathrm{cont}} = U_{\mathrm{American}} - \Phi.
\]
American-minus-European value is:
\[
U_{\mathrm{American}} - U_{\mathrm{European}}.
\]

The first quantity compares the American value to immediate exercise. It is the right quantity
for obstacle and boundary extraction. The second compares the American contract to a European
contract. It is useful for validation, especially in Tickets 05 through 07, but it is not the
boundary-extraction object.

\section{Implementation Design Decisions}

\begin{longtable}{L{0.30\textwidth}L{0.58\textwidth}}
\toprule
Decision & Reason \\
\midrule
\endhead
Create a diagnostics module & Boundary extraction is a diagnostic layer on solver outputs, not part of the solver core. \\
Use Ticket 04 result objects & The stored value grid, payoff, spot grid, and tau grid already contain the needed data. \\
Keep solver files unchanged & Ticket 09 should not introduce a new pricing algorithm or alter CN/PSOR logic. \\
Use full grids for profiles & Full profiles help students see the premium shape across the domain. \\
Use interior nodes for extraction & Interior extraction avoids treating imposed boundary values as solved boundary locations. \\
Record no-boundary reasons & Missing boundaries are informative, especially for the no-dividend call control. \\
Include no-dividend control & It checks that the extractor does not force a high-spot call boundary when \(q=0\). \\
\bottomrule
\end{longtable}

\section{Code-Level Function Explanations}

\begin{longtable}{P{0.35\textwidth}L{0.53\textwidth}}
\toprule
Function or class & Role \\
\midrule
\endhead
\func{BoundaryPoint} & Stores one time-level boundary estimate or no-boundary reason. \\
\func{BoundaryCurve} & Stores all time-level boundary points plus value, payoff, and premium grids. \\
\func{BoundaryExtractionSummary} & Stores whole-curve counts, boundary ranges, and status. \\
\func{continuation_premium} & Computes \(U-\Phi\) for one row or a full value grid. \\
\func{linear_interpolate_threshold_crossing} & Estimates the threshold-crossing spot between adjacent nodes. \\
\func{extract_boundary_at_time} & Applies the option-type search rule to one premium row. \\
\func{extract_boundary_curve} & Applies boundary extraction over all time-to-maturity rows. \\
\func{summarize_boundary_curve} & Converts all boundary points into summary counts and ranges. \\
\func{selected_time_profile_rows} & Exports full-grid premium rows at selected time fractions for study and plotting. \\
\func{run_boundary_cases} & Runs the three Ticket 09 experiment cases and prepares CSV/figure artifacts. \\
\bottomrule
\end{longtable}

\section{Extraction Cases and Parameters}

\begin{longtable}{L{0.24\textwidth}L{0.10\textwidth}L{0.08\textwidth}L{0.08\textwidth}L{0.08\textwidth}L{0.08\textwidth}L{0.09\textwidth}L{0.08\textwidth}L{0.08\textwidth}}
\toprule
Case & Option & \(K\) & \(T\) & \(r\) & \(q\) & \(\sigma\) & \(M\) & \(N\) \\
\midrule
\endhead
Put medium & put & 1.0 & 1.0 & 0.05 & 0.02 & 0.20 & 120 & 120 \\
Dividend call & call & 1.0 & 1.0 & 0.05 & 0.08 & 0.20 & 120 & 120 \\
No-dividend call & call & 1.0 & 1.0 & 0.05 & 0.00 & 0.20 & 120 & 120 \\
\bottomrule
\end{longtable}

All cases use \(S_{\max}=4.0\) and threshold \(10^{-6}\). The no-dividend call control is not a
repeat of Ticket 06 price validation. It is a boundary-extraction control.

\section{Boundary Result Tables and Interpretation}

\subsection{Boundary Summary}

\begin{longtable}{L{0.25\textwidth}L{0.12\textwidth}L{0.13\textwidth}L{0.13\textwidth}L{0.12\textwidth}L{0.12\textwidth}}
\toprule
Case & Direction & Found rows & No-boundary rows & Min boundary & Max boundary \\
\midrule
\endhead
Put medium & low-to-high & 120 & 1 & 0.766754 & 0.934082 \\
Dividend call & high-to-low & 120 & 1 & 1.066165 & 1.266630 \\
No-dividend call & high-to-low & 0 & 121 & n/a & n/a \\
\bottomrule
\end{longtable}

The put and dividend-call cases find boundaries for every positive-\(\tau\) row. The maturity row
is marked ambiguous because \(U=\Phi\) everywhere at \(\tau=0\). The no-dividend call control finds
no boundary, which is the desired behavior for a control case.

\subsection{No-Boundary Reason Counts}

\begin{longtable}{L{0.28\textwidth}L{0.12\textwidth}L{0.12\textwidth}L{0.18\textwidth}L{0.14\textwidth}}
\toprule
Case & Maturity & All cont. & Expected side absent & No clean transition \\
\midrule
\endhead
Put medium & 1 & 0 & 0 & 0 \\
Dividend call & 1 & 0 & 0 & 0 \\
No-dividend call & 1 & 0 & 120 & 0 \\
\bottomrule
\end{longtable}

The control case reports \path{expected_exercise_side_absent} for positive time rows. In plain
language, the high-\(S\) side of the no-dividend call does not show an exercise-like transition.

\subsection{Selected By-Time Rows}

\begin{longtable}{L{0.25\textwidth}L{0.07\textwidth}L{0.09\textwidth}L{0.12\textwidth}L{0.16\textwidth}L{0.14\textwidth}}
\toprule
Case & Step & \(\tau\) & Found & Boundary spot & Reason \\
\midrule
\endhead
put & 0 & 0.00000 & false & n/a & maturity ambiguous \\
put & 1 & 0.00833 & true & 0.934082 & -- \\
put & 60 & 0.50000 & true & 0.800004 & -- \\
put & 120 & 1.00000 & true & 0.766754 & -- \\
dividend call & 0 & 0.00000 & false & n/a & maturity ambiguous \\
dividend call & 1 & 0.00833 & true & 1.066165 & -- \\
dividend call & 60 & 0.50000 & true & 1.200001 & -- \\
dividend call & 120 & 1.00000 & true & 1.266630 & -- \\
no-dividend call & 1 & 0.00833 & false & n/a & expected side absent \\
no-dividend call & 120 & 1.00000 & false & n/a & expected side absent \\
\bottomrule
\end{longtable}

These rows are selected for readability. The full by-time CSV contains 363 rows: three cases times
121 time rows.

\section{Figure Interpretation}

Figure~\ref{fig:put-boundary} plots the American put boundary over time-to-maturity. The boundary
is below the strike and moves gradually on the tested grid.

\begin{figure}[H]
\centering
\includegraphics[width=0.82\textwidth]{results/01_solver_validation/figures/ticket_09_american_put_boundary_curve.png}
\caption{Ticket 09 American put boundary curve extracted from continuation premium.}
\label{fig:put-boundary}
\end{figure}

Figure~\ref{fig:call-boundary} plots the dividend-paying American call boundary. It appears above
the strike in this representative \(q=0.08\) case.

\begin{figure}[H]
\centering
\includegraphics[width=0.82\textwidth]{results/01_solver_validation/figures/ticket_09_dividend_call_boundary_curve.png}
\caption{Ticket 09 dividend-paying American call boundary curve extracted from continuation premium.}
\label{fig:call-boundary}
\end{figure}

Figure~\ref{fig:premium-profiles} shows selected premium profiles. These curves help students see
why the boundary is a threshold-crossing diagnostic rather than a visual guess.

\begin{figure}[H]
\centering
\includegraphics[width=0.92\textwidth]{results/01_solver_validation/figures/ticket_09_continuation_premium_profiles.png}
\caption{Ticket 09 continuation-premium profiles at selected time-to-maturity levels.}
\label{fig:premium-profiles}
\end{figure}

Figure~\ref{fig:status} records whether a boundary was found at each time row. The no-dividend
call control stays at ``not found'' for positive time rows.

\begin{figure}[H]
\centering
\includegraphics[width=0.86\textwidth]{results/01_solver_validation/figures/ticket_09_boundary_found_status.png}
\caption{Ticket 09 boundary-found status by case and time-to-maturity row.}
\label{fig:status}
\end{figure}

\section{Testing Strategy}

\begin{longtable}{L{0.30\textwidth}L{0.30\textwidth}L{0.28\textwidth}}
\toprule
Test category & What it proves & What it does not prove \\
\midrule
\endhead
Premium calculation & \(U-\Phi\) is computed with correct broadcasting. & Boundary accuracy. \\
Synthetic put boundary & Low-spot put search direction works on a controlled transition. & Real put convergence. \\
Synthetic call boundary & High-spot dividend-call search direction works on a controlled transition. & Real call convergence. \\
No-boundary metadata & Missing boundaries are explained rather than hidden. & That every future missing boundary is harmless. \\
Interpolation & The threshold crossing formula is predictable. & Higher-order boundary accuracy. \\
Real put and call smoke cases & Existing solver outputs produce boundary rows and nonnegative premiums. & Full grid/domain sensitivity. \\
No-dividend control & The extractor does not force a high-spot call boundary for \(q=0\). & Ticket 06 theorem validation. \\
CSV and figure artifacts & Reproducible tables and useful figures are produced. & Scientific proof by plotting. \\
Scope and docstrings & Ticket 09 files identify the ticket and avoid Greeks, stress maps, datasets, and neural work. & Later-ticket functionality. \\
\bottomrule
\end{longtable}

\section{Test Results and Interpretation}

The Ticket 09 implementation followed a test-first workflow. The first red run failed because
\file{src/american_risk_surfaces/diagnostics/boundary.py} and
\file{experiments/06_boundary_extraction.py} did not exist yet. After implementation, the targeted
Ticket 09 test command reported:

\begin{lstlisting}[language=bash]
Ran 17 tests in 0.611s
OK
\end{lstlisting}

The full test command reported:
\begin{lstlisting}[language=bash]
Ran 104 tests in 1.824s
OK
\end{lstlisting}

This means the Ticket 09 tests passed together with all earlier Ticket 01--08 tests. It does not
mean the extracted boundaries are exact analytical free boundaries.

\section{Implementation Pitfalls and Debugging Notes}

\begin{longtable}{L{0.25\textwidth}L{0.25\textwidth}L{0.22\textwidth}L{0.16\textwidth}}
\toprule
Recorded item or pitfall & Symptom & Fix or guard & Beginner lesson \\
\midrule
\endhead
Expected TDD red run & New tests failed because the Ticket 09 module and experiment did not exist. & Add the planned module and experiment. & Red tests confirm new behavior is actually being tested. \\
Malformed targeted unittest command & The dotted test-module form could not import because \file{tests/} is not a package. & Use unittest discovery with the tests directory and filename pattern. & Match test commands to repo layout. \\
Overly strict transition rule & Real rows looked like exercise-like, continuation-like, exercise-like because far tails can have tiny premium. & Use option-type search direction and require a minimum continuation block near the expected exercise side. & Boundary extraction needs financial direction, not only boolean pattern matching. \\
Confusing premium types & American-minus-European value is tempting to plot. & Use \(U-\Phi\) for boundary extraction. & Boundary extraction uses payoff comparison, not European comparison. \\
Forcing a control boundary & A no-dividend call may have tiny low-spot premium. & Search high-to-low for calls and report expected-side absence for \(q=0\). & Negative controls are valuable. \\
Overclaiming exactness & A smooth-looking curve can feel authoritative. & Report threshold, interpolation, and limitations. & Numerical boundaries are estimates with metadata. \\
\bottomrule
\end{longtable}

Future coding/test reports should continue to include implementation pitfalls when failed tests,
debugging discoveries, or workflow corrections occur.

\section{Limitations}

\begin{longtable}{L{0.34\textwidth}L{0.54\textwidth}}
\toprule
Not included in Ticket 09 & Reason \\
\midrule
\endhead
New pricing algorithm & Ticket 09 reads existing American CN/PSOR outputs only. \\
Greeks & Delta and Gamma belong to Ticket 10. \\
Greek masks near boundary & Masking depends on later Greek diagnostics. \\
Stress maps & Regime sweeps are later validation work. \\
Datasets and neural networks & Labels and models require a more mature validated solver pipeline. \\
Exact free-boundary accuracy claim & The current output is threshold-based and grid-dependent. \\
Full grid/domain sensitivity & This ticket uses representative \(M=N=120\) cases only. \\
\bottomrule
\end{longtable}

\section{Lessons Learned / Study Notes}

The first lesson is that a free boundary is not drawn by eye. It should come from a defined
quantity, here the continuation premium \(U-\Phi\).

The second lesson is that search direction matters. Put boundaries and dividend-call boundaries
live on different sides of the spot domain.

The third lesson is that missing boundaries are data, not failure by default. The no-dividend
American call control should not have a forced high-spot exercise boundary.

The fourth lesson is that thresholds must be reported. A boundary extracted with \(\epsilon=10^{-6}\)
is not the same claim as a boundary extracted with a much larger threshold.

The fifth lesson is to keep later tickets separate. Boundary extraction is now allowed, but Greeks,
stress maps, datasets, and neural-network labels are still out of scope.

\section{Link to Ticket 10}

Ticket 10 should compute Delta and Gamma diagnostics. It can use the Ticket 09 boundary curves to
define careful reporting bands near the free boundary, but it should not treat the current
boundary estimates as exact analytical truth. Greek diagnostics should remain honest about payoff
kinks, threshold effects, and boundary-near numerical instability.

\section{Reproducibility Record}

\begin{longtable}{L{0.30\textwidth}L{0.58\textwidth}}
\toprule
Item & Record \\
\midrule
\endhead
Boundary module & \file{src/american_risk_surfaces/diagnostics/boundary.py}. \\
Experiment & \file{experiments/06_boundary_extraction.py}. \\
Tests & \file{tests/test_boundary_extraction.py}. \\
Summary CSV & \file{results/01_solver_validation/tables/ticket_09_boundary_extraction_summary.csv}. \\
By-time CSV & \file{results/01_solver_validation/tables/ticket_09_boundary_by_time.csv}. \\
Selected-times CSV & \file{results/01_solver_validation/tables/ticket_09_boundary_selected_times.csv}. \\
Figures & \file{results/01_solver_validation/figures/ticket_09_american_put_boundary_curve.png};
\file{results/01_solver_validation/figures/ticket_09_dividend_call_boundary_curve.png};
\file{results/01_solver_validation/figures/ticket_09_continuation_premium_profiles.png};
\file{results/01_solver_validation/figures/ticket_09_boundary_found_status.png}. \\
Report source & \file{reports/03_solver/tickets/ticket_09_continuation_premium_boundary_extraction.tex}. \\
Report PDF & \file{reports/03_solver/tickets/ticket_09_continuation_premium_boundary_extraction.pdf}. \\
\bottomrule
\end{longtable}

Commands used:
\begin{lstlisting}[language=bash]
PYTHONPYCACHEPREFIX=/private/tmp/codex_pycache PYTHONPATH=src MPLCONFIGDIR=/private/tmp/matplotlib-cache FC_CACHEDIR=/private/tmp/fontconfig-cache python3 experiments/06_boundary_extraction.py
PYTHONPYCACHEPREFIX=/private/tmp/codex_pycache PYTHONPATH=src MPLCONFIGDIR=/private/tmp/matplotlib-cache FC_CACHEDIR=/private/tmp/fontconfig-cache python3 -m unittest discover -s tests
PYTHONPYCACHEPREFIX=/private/tmp/codex_pycache PYTHONPATH=src python3 -m compileall -q src tests experiments
\end{lstlisting}

\bibliographystyle{plainnat}
\bibliography{reports/03_solver/references}

\end{document}
